\section{2026-05-18}

\sessionhead{1}{14:02}{spike}{decoder-explorer}{(none)}

\subsubsection*{Context --- what this experiment is about}
The Tunable-ROM paper uses a \textbf{CP decoder} inside its non-linear
manifold reduced-order model (NM-ROM). CP works well, but we wanted to
ask a wider question: \emph{can a different kind of autoencoder do the
same job and maybe be faster?} Concretely, can we swap the CP decoder
for a \textbf{SIREN} decoder (coordinate-based MLP with sine
activations, a popular ``implicit neural representation'' or INR), keep
EQ hyper-reduction (the trick that makes the ROM cheap), and still
solve the PDE through the ROM?

If yes, we get a more flexible decoder family --- INRs can be queried
at any spatial point, not just grid nodes, which is useful for
refinement, multi-resolution, and irregular domains. If no, we learn
\emph{why} (Phase 4 in \texttt{JOURNAL.tex} is that negative result for
the cross-attention INR variant). This session is the SIREN follow-up.

\paragraph{Setup.} 2-D Poisson on the unit square, zero-Dirichlet BC,
parametric forcing $f(x;k) = A \prod_i \sin(k_i \pi x_i)$,
$k_i \in [1,3]$, $A=10$. Grid $N{=}128$ (16{,}384 unknowns). FOM is a
sparse 5-point Laplacian with Conjugate Gradient (median $\sim$270 ms
per parameter). Data: 700 train / 140 val / 160 test snapshots all
generated by that same CG FOM.

\paragraph{Model.} ViT encoder (patch=16, embed=64, 4 layers,
4 heads) $\rightarrow$ 16-dim latent $z$ $\rightarrow$ modulated SIREN
decoder (5 sine layers, hidden=384, $\omega_0{=}30$, modulator MLP
mapping $z$ to per-layer bias offsets). The decoder is queried
\emph{point by point}: give it a coord $x$, it returns $u(x)$. That
point-wise query is the whole appeal of an INR.

\paragraph{Solve.} For a new parameter $\mu$, pick a starting latent
$z_0$, then run Levenberg--Marquardt Gauss--Newton on $z$ to drive the
discrete Poisson residual at the EQ stencil nodes to zero. Two natural
choices for $z_0$: \emph{warm-start} ($z_0 = \text{encoder}(\text{
analytical}\_u(\mu))$, requires an oracle solution) or \emph{cold-start}
($z_0 = 0$, no oracle needed). \textbf{The whole session is about making
cold-start work}, because it's the production-relevant case.

\paragraph{Starting state.} Pre-session SIREN\_WIDE checkpoint trained
cleanly to AE val rel-L$^2 = 8.88$e-3 --- the autoencoder works. But:
\textbf{cold-start ROM solve rel-L$^2$ $= 1.565$ for every test
parameter, every $n_{\text{eq}}$}; warm-start works at rel-L$^2 =
5.12$e-2 ($\sim$16$\times$ worse than the CP-decoder reference of
3.23e-3). Target: close that gap from cold-start.

\subsubsection*{Summary}
Cracked the cold-start NM-ROM failure for SIREN on Poisson-2D, $N{=}128$:
median rel-L$^2$ dropped from 1.565 to 2.88e-3 (matches the CP-decoder
reference) across 7 iterations of training-recipe and solver-side changes.
A diagnostic showed the failure was a spurious basin of attraction at
$z{=}0$, not a Jacobian-rank issue.

\subsubsection*{Files touched}
\begin{itemize}
  \item \texttt{Experiments/exploring-decoders/scripts/diagnose\_siren\_cold.py}
        --- new diagnostic: SVD of $J(z{=}0)$, residual--range alignment,
        1-D landscape, and 7 alternative $z_0$ tests on a single $\mu$.
  \item \texttt{Experiments/exploring-decoders/src/tunable\_rom\_decoders/training\_rom\_anchor.py}
        --- anchor-at-zero training (force $\text{decode}(0,\cdot) \approx
        \text{mean}(U_{\text{train}})$, plus $\lambda_z \cdot \|z\|^2$
        compression). Used by iter 2 (the breakthrough).
  \item \texttt{Experiments/exploring-decoders/src/tunable\_rom\_decoders/training\_rom\_lap.py}
        --- anchor plus ROM-aware Laplacian residual. Used by iter 5/6.
  \item \texttt{Experiments/exploring-decoders/scripts/run\_rom\_continuation.py}
        --- two-stage solver (coarse $n_{\text{eq}}$ $\rightarrow$
        warm-start fine $n_{\text{eq}}$) on a single checkpoint.
  \item \texttt{Experiments/exploring-decoders/scripts/run\_rom\_composite.py}
        --- composite two-checkpoint solver: iter 2 cold-seed $\rightarrow$
        field $\rightarrow$ iter 5 encode $\rightarrow$ iter 5 warm-refine.
        Produced the final SOTA cold-start result.
  \item \texttt{Experiments/exploring-decoders/src/tunable\_rom\_decoders/solver/nm\_rom.py}
        --- added \texttt{encode\_method\_name},
        \texttt{latent\_dim\_override}, and \texttt{gn\_max\_iters} plumbing
        to support VDF and cross-checkpoint use.
  \item \texttt{Experiments/exploring-decoders/journey-to-good-manifold.md}
        --- per-iteration log with a ``Mechanistic-interpretability
        post-mortem of Iteration 0'' section walking through the four GN
        failure modes and which probe killed each.
  \item \texttt{Experiments/exploring-decoders/JOURNAL.tex} --- added
        \S Phase 5 summarizing iters 0--7.
  \item Also created: \texttt{training\_rom\_vdf.py},
        \texttt{training\_rom\_zreg.py}, \texttt{autoencoder\_vdf.py},
        \texttt{default\_field\_siren.py}, \texttt{linear\_skip\_siren.py},
        plus several configs and SLURM scripts (not all load-bearing for
        the final result).
\end{itemize}

\subsubsection*{Decisions / surprises}
\begin{itemize}
  \item \textbf{Failure mode was C (spurious basin), not A/B/D.} The
        diagnostic measured $\text{cond}\,J(z{=}0) = 17.2$ (full rank,
        kills A) and $\|\text{proj}_{\text{range}\,J}\,R\| / \|R\| = 0.69$
        (most residual energy is reachable, kills B). The 1-D landscape
        showed good basins within $\|z\| \approx 1$--$2$ of zero, so the
        iter cap (D) was not the limiter. The smoking gun: 7 different
        $z_0$ seeds for one test $\mu$ --- every ``small'' seed (zero,
        randn $\sigma{=}0.01$, randn $\sigma{=}0.1$) converged to
        $\|z_{\text{final}}\| \approx 0.38$ with rel-L$^2 \approx 1.3$,
        while $\text{encoder}(\text{analytical}\_u(\mu))$ at
        $\|z_0\|{=}1.57$ converged in 7 iters to rel-L$^2 = 0.043$. There
        is a well-defined bad attractor near the origin; the manifold's
        image of $z{=}0$ is just out of distribution.
  \item \textbf{Two scalar loss terms fix it.} $\lambda_z \cdot \|z\|^2$
        on the encoder output compresses trained latents from
        $\|z\| \approx 1.5$ to $\|z\| \approx 0.14$.
        $\lambda_a \cdot \text{rel\_l2}(\text{decode}(0,\cdot),
        \text{mean}(U_{\text{train}})(\cdot))$ makes $z{=}0$ decode to
        the mean training field. With $\lambda_z{=}0.1$ and
        $\lambda_a{=}1.0$ (iter 2), cold-start median rel-L$^2$ goes
        $1.565 \rightarrow 0.043$ with no architecture change and no
        extra inference cost.
  \item \textbf{The recipe is fragile.} Scaling SIREN width
        ($h{=}384 \rightarrow 512$, iter 4) keeps warm-start working but
        completely breaks cold-start (median 0.945). Adding a ROM-aware
        Laplacian loss at any weight (iter 5, $\lambda_L{=}0.1$; iter 6,
        $\lambda_L{=}0.01$) also breaks cold-start --- but lifts
        warm-start to median 2.98e-3 in 7 iters (matches the CP-decoder
        reference of 3.23e-3). The two training objectives (anchor
        $\text{decode}(0) = u_{\text{mean}}$ vs.\ $K \cdot
        \text{decode}(z) \approx F_i$ for all $i$) are structurally
        inconsistent.
  \item \textbf{Resolution: compose at solve time.} Iter 7 uses iter 2
        to cold-start at $n_{\text{eq}}{=}500$ (gets a decent $z$),
        decodes to a full-grid field, re-encodes with iter 5's encoder,
        then warm-starts iter 5's solver. Final cold-start median
        rel-L$^2 = 2.88$e-3, 94\% of test parameters below 1e-2, max
        1.42 (a $\sim$6\% high-frequency tail where iter 2's seed is
        poor). Cost: 2 solves $+$ 1 encode per $\mu$ --- speedup over
        FOM is only 1.24$\times$ at fine $n_{\text{eq}}{=}500$,
        0.5$\times$ at fine $n_{\text{eq}}{=}8000$.
  \item \textbf{Honest accuracy-vs-speed split.} This session's win is
        accuracy parity with the CP-decoder reference (matching the
        paper number from cold start $z{=}0$). It is NOT a speed win ---
        INRs decode point-by-point so every GN iteration costs $\sim$30
        GFLOPs (vs $\sim$3 GFLOPs for the sparse FOM matvec), and EQ
        hyper-reduction can't paper over that gap. This is consistent
        with the Phase 4 conclusion already in \texttt{JOURNAL.tex}.
\end{itemize}

\subsubsection*{Results / metrics}

Cold-start ($z_0 = 0$) on the 160-sample test set, $n_{\text{eq}}{=}500$
(or coarse 500 / fine 500 for iter 7), \texttt{gn\_max\_iters} $=30$ for
iter 3+, 12 for iter 0/1b/2:

\begin{center}
\begin{tabular}{l r r r r r}
  \textbf{Recipe} & \textbf{median rel-L$^2$} & \textbf{p90 rel-L$^2$} & \textbf{AE val rel-L$^2$} & \textbf{iters (med.)} & \textbf{speedup} \\
  \hline
  baseline SIREN\_WIDE                                & 1.565   & ---     & 8.88e-3 & 12 & 1.0$\times$ \\
  iter 1b: $+$ L2($z$) $\lambda{=}$1e-3               & 1.311   & ---     & 4.01e-3 & 12 & 1.2$\times$ \\
  iter 2: $+$ anchor $+$ L2($z$) $\lambda{=}$1e-1     & 4.25e-2 & 6.8e-1  & 7.85e-3 & 12 & 5.1$\times$ \\
  iter 3a: iter2 ckpt $+$ 30 GN iters                 & 2.33e-2 & 2.4e-1  & (same as iter2) & 30 & 2.1$\times$ \\
  iter 3b: iter2 ckpt $+$ EQ continuation (500$\to$8000) & 2.30e-2 & 1.0e-1 & (same as iter2) & 30 & 0.1$\times$ \\
  iter 4: iter2 recipe $+$ XL width ($h{=}512$)       & 0.945   & 1.07    & 1.09e-2 & 30 & 1.4$\times$ \\
  iter 5: iter2 $+$ Laplacian $\lambda_L{=}$1e-1      & 0.918   & 1.24    & 1.10e-2 & 30 & 2.1$\times$ \\
  iter 6: iter2 $+$ Laplacian $\lambda_L{=}$1e-2      & 0.909   & 0.97    & 1.03e-2 & 30 & 2.1$\times$ \\
  \textbf{iter 7: composite iter2 $\to$ iter5 (500$\to$500)} & \textbf{2.88e-3} & \textbf{6.6e-3} & (composite) & \textbf{5 (fine)} & \textbf{1.24$\times$} \\
\end{tabular}
\end{center}

Warm-start ($z_0 = \text{encoder}(\text{analytical}\_u(\mu))$, one test
parameter, $n_{\text{eq}}{=}500$, 30 GN iters cap):

\begin{center}
\begin{tabular}{l r r r}
  \textbf{Recipe} & \textbf{median rel-L$^2$} & \textbf{iters (med.)} & \textbf{speedup} \\
  \hline
  baseline SIREN\_WIDE                            & 5.12e-2 & 12 & 1.0$\times$ \\
  iter 2: anchor $+$ L2($z$)                      & 1.84e-2 & 12 & 5.1$\times$ \\
  iter 4: XL width                                & 1.90e-2 & 14 & 2.8$\times$ \\
  iter 6: Laplacian $\lambda_L{=}$1e-2            & 8.10e-3 & 6  & 9.4$\times$ \\
  \textbf{iter 5: Laplacian $\lambda_L{=}$1e-1}   & \textbf{2.98e-3} & \textbf{7} & \textbf{8.4$\times$} \\
\end{tabular}
\end{center}

\sessionhead{2}{14:51}{spike}{inr-siren-speed-accuracy}{b3c9325, 3fb4c16, 312ee5d, a3a5320, 474359e, 1ab8ddd}

\subsubsection*{Summary}
Pushed cold-start NM-ROM from 0.13$\times$ FOM at rel-L$^2$ 2.88e-3 to a
three-tier Pareto curve spanning 3.4$\times$--320$\times$ FOM by replacing
\texttt{jax.jacfwd} with \texttt{vmap(jacrev)} and introducing a
linear-affine decoder $u(x;z) = V(x){\cdot}z + b(x)$ trained with a
Laplacian residual loss.

\subsubsection*{Files touched}
\begin{itemize}
  \item \texttt{Experiments/2026-05-18-INR-SIREN-SPEED-ACCURACY-TRADEOFF/src/tunable\_rom\_speed/solver/nm\_rom\_fast.py}
        --- \texttt{FastINRNMROMSolver}: per-iter Jacobian via
        \texttt{vmap(jacrev)} instead of \texttt{jacfwd}.
  \item \texttt{Experiments/2026-05-18-INR-SIREN-SPEED-ACCURACY-TRADEOFF/src/tunable\_rom\_speed/decoders/linear\_affine\_z.py}
        --- \texttt{LinearAffineZDecoder}: $u(x;z) = V(x){\cdot}z + b(x)$
        with constant Jacobian in $z$.
  \item \texttt{Experiments/2026-05-18-INR-SIREN-SPEED-ACCURACY-TRADEOFF/src/tunable\_rom\_speed/solver/nm\_rom\_affine.py}
        --- \texttt{AffineNMROMSolver}: precomputes $V(x_{\text{eq}})$
        once per solve.
  \item \texttt{Experiments/2026-05-18-INR-SIREN-SPEED-ACCURACY-TRADEOFF/scripts/timing\_harness.py}
        --- per-phase millisecond decomposition.
  \item \texttt{Experiments/2026-05-18-INR-SIREN-SPEED-ACCURACY-TRADEOFF/scripts/diagnose\_affine\_manifold.py}
        --- oracle-vs-ROM rel-L$^2$ split for the affine decoder.
\end{itemize}

\subsubsection*{Decisions / surprises}
\begin{itemize}
  \item \textbf{Iter-7 baseline was 0.13$\times$ FOM, not 0.50$\times$.}
        Prior timing was JIT-warmed and undercounted; the honest cold
        wall is 2152 ms.
  \item \textbf{Nonlinear $A(z)$ in the affine decoder failed
        (rel-L$^2$ 0.2--0.8)} because cold-start GN got trapped in a
        basin near $z{=}0$. Switching to linear $A(z)$ (constant
        Jacobian) fixed it.
  \item \textbf{LinearAffine alone solved the manifold but left an
        encoder-vs-ROM consistency gap} (oracle 1e-2 but ROM 0.25).
        Adding a Laplacian residual loss during training closed it
        ($\|z_{\text{enc}}\| \approx \|z_{\text{rom}}\|$ within 1\%).
  \item \textbf{Iter-7's $n_{\text{eq}}{=}8000$ was load-bearing for the
        bad speedup.} $n_{\text{eq}}{=}500$ with 10--20 GN iters gives
        identical accuracy at $\sim$1/8 the per-iter cost.
\end{itemize}

\subsubsection*{Results / metrics}

Cold-start Pareto curve on the 160-sample test set:

\begin{center}
\begin{tabular}{l l r r r r}
  \textbf{Tier} & \textbf{Decoder} & \textbf{median rel-L$^2$} & \textbf{p90} & \textbf{speedup vs FOM} & \textbf{wall (ms)} \\
  \hline
  iter-7 baseline (re-timed) & SIREN composite       & 2.88e-3 & 6.61e-3 & 0.13$\times$ & 2152 \\
  Accurate                   & Fast SIREN ic30/if20  & 2.88e-3 & 6.66e-3 & 3.44$\times$ & 80   \\
  Balanced                   & Fast SIREN ic20/if10  & 2.89e-3 & 9.87e-3 & 3.78$\times$ & 72   \\
  Sub-1\%                    & linaff\_v5            & 9.81e-3 & 1.42e-2 & 139$\times$  & 1.97 \\
  Speed-first                & linaff\_v2            & 1.81e-2 & 3.16e-2 & 320$\times$  & 0.85 \\
\end{tabular}
\end{center}

LinAff training-curve scaling (winner in bold):

\begin{center}
\begin{tabular}{l l r r r r r}
  \textbf{Config} & \textbf{$\Phi$} & \textbf{$z$-dim} & \textbf{epochs} & \textbf{$\lambda_{\text{lap}}$} & \textbf{val rel-L$^2$} & \textbf{ROM rel-L$^2$} \\
  \hline
  linaff\_v1 & 384w/4L & 16 & 50k  & 0     & 1.0e-2  & 0.25    \\
  linaff\_v2 & 512w/5L & 16 & 50k  & 0.1   & 5.5e-2  & 1.81e-2 \\
  linaff\_v4 & 512w/5L & 16 & 100k & 0.2   & 8.6e-3  & 1.31e-2 \\
  linaff\_v3 & 640w/6L & 16 & 100k & 0.05  & 9.6e-3  & 1.14e-2 \\
  \textbf{linaff\_v5 (winner)} & \textbf{768w/5L} & \textbf{24} & \textbf{150k} & \textbf{0.15} & \textbf{4.2e-3} & \textbf{9.81e-3} \\
\end{tabular}
\end{center}

\subsubsection*{Architecture appendix}

This appendix explains the four modules built today in enough detail that
the design choices make sense without re-reading the source. Read it as
two pairs: a solver-side speedup (\texttt{FastINRNMROMSolver}) on top of
the inherited decoder (\texttt{ModulatedSIREN}), and a decoder-side
speedup (\texttt{LinearAffineZDecoder}) reached by first failing with
\texttt{AffineZDecoder}.

\begin{figure}[H]
  \centering
  \begin{tikzpicture}[
    >={Stealth[length=2mm]},
    node distance=4mm and 6mm,
    box/.style={draw, rounded corners=2pt, minimum height=7mm, minimum width=14mm, align=center, font=\small},
    fixed/.style={box, fill=blue!8},
    learned/.style={box, fill=orange!12},
    output/.style={box, fill=green!10},
    zlatent/.style={circle, draw, fill=red!12, minimum size=6mm, inner sep=0pt, font=\small},
    flow/.style={->, thick},
  ]
    % ===== Column 1: ModulatedSIREN =====
    \node[font=\small\bfseries] (t1) at (0,0) {ModulatedSIREN};
    \node[zlatent, below=4mm of t1] (z1) {$z$};
    \node[learned, below=5mm of z1] (mod1) {Modulator\\MLP};
    \node[learned, below=4mm of mod1] (trunk1) {SIREN trunk\\$\sin(W h{+}b{+}\beta)$};
    \node[font=\small, left=5mm of trunk1] (x1) {$x$};
    \node[output, below=4mm of trunk1] (u1) {$u(x)$};
    \draw[flow] (z1) -- (mod1);
    \draw[flow] (mod1) -- node[right, font=\scriptsize] {$\beta(z)$} (trunk1);
    \draw[flow] (x1) -- (trunk1);
    \draw[flow] (trunk1) -- (u1);
    \node[below=2mm of u1, font=\scriptsize, align=center, text width=42mm] (c1)
      {$\partial u/\partial z$: nonlinear, dense.\\
       Per-iter: jacfwd ($d_z$ fwds)\\
       or vmap(jacrev) (1 fwd $+$ 1 bwd / pt).};

    % ===== Column 2: AffineZDecoder =====
    \node[font=\small\bfseries, right=22mm of t1] (t2) {AffineZDecoder (failed)};
    \node[zlatent, below=4mm of t2] (z2) {$z$};
    \node[learned, below=5mm of z2] (Anl) {Nonlinear\\$A(z)$ MLP};
    \node[font=\small, left=14mm of Anl] (x2) {$x$};
    \node[learned] (phi2) at (x2 |- Anl) {$\Phi(x)$\\trunk};
    \node[learned, below=4mm of phi2] (b2) {$b(x)$\\trunk};
    \node[output] (u2) at ($(Anl)!0.5!(b2) + (0,-10mm)$) {$u(x){=}\Phi(x){\cdot}A(z){+}b(x)$};
    \draw[flow] (z2) -- (Anl);
    \draw[flow] (x2) -- (phi2);
    \draw[flow] (x2.south) |- (b2.west);
    \draw[flow] (phi2) -- (u2);
    \draw[flow] (Anl) -- (u2);
    \draw[flow] (b2) -- (u2);
    \node[below=2mm of u2, font=\scriptsize, align=center, text width=46mm] (c2)
      {$\partial u/\partial z = \Phi(x)\cdot \partial A/\partial z$,\\
       non-constant in $z$ --- basin near $z{=}0$.};

    % ===== Column 3: LinearAffineZDecoder =====
    \node[font=\small\bfseries, right=26mm of t2] (t3) {LinearAffineZDecoder (fix)};
    \node[zlatent, below=4mm of t3] (z3) {$z$};
    \node[fixed, below=5mm of z3] (Alin) {Linear\\$A(z){=}W_A\,z$};
    \node[font=\small, left=14mm of Alin] (x3) {$x$};
    \node[learned] (phi3) at (x3 |- Alin) {$\Phi(x)$\\trunk};
    \node[learned, below=4mm of phi3] (b3) {$b(x)$\\trunk};
    \node[fixed] (V3) at ($(phi3)!0.5!(Alin) + (0,-10mm)$) {$V(x){=}\Phi(x)\,W_A$\\(precomputed)};
    \node[output, below=4mm of V3] (u3) {$u(x){=}V(x){\cdot}z{+}b(x)$};
    \draw[flow] (z3) -- (Alin);
    \draw[flow] (x3) -- (phi3);
    \draw[flow] (x3.south) |- (b3.west);
    \draw[flow] (phi3) -- (V3);
    \draw[flow] (Alin) -- (V3);
    \draw[flow] (V3) -- (u3);
    \draw[flow] (b3.east) -| (u3.north);
    \node[below=2mm of u3, font=\scriptsize, align=center, text width=46mm] (c3)
      {$\partial u/\partial z = V(x)$, constant in $z$\\
       --- GN converges in 1--2 iters.};
  \end{tikzpicture}
  \caption*{\small Three decoder architectures contrasted. Red = latent $z$.
  Orange = MLPs whose output depends on $z$ (must be re-evaluated each GN
  iter). Blue = computations independent of $z$ (precomputable once per
  solve). Green = decoded scalar $u(x)$.
  The Jacobian $\partial u/\partial z$ rank/structure is annotated below each.}
\end{figure}

\paragraph{\texttt{FastINRNMROMSolver}.}
A drop-in subclass of \texttt{INRNMROMSolver} that changes only how the
Gauss--Newton Jacobian is built. The baseline used
\texttt{jax.jacfwd(residual\_vec)(z)}, which runs \emph{latent\_dim}
parallel forward passes of the full $(n_{\text{eq}},)$-output residual
to fill one column per latent component. The fast version pushes the
Jacobian inside the per-point loop: each EQ stencil point has a scalar
decoder output, so \texttt{jacrev} costs one forward $+$ one backward
per point, and \texttt{vmap} reuses the kernel across the
$n_{\text{eq}}$ points. Math is unchanged --- the finite-difference
stencil combination after the per-point Jacobians is identical --- and
correctness was checked to $\|z_{\text{fast}}-z_{\text{base}}\| /
\|z_{\text{base}}\| \le 6.7$e-9.

\[
J_{\text{base}} \;=\; \mathtt{jacfwd}(r)(z)\,, \qquad
J_{\text{fast}}[i,:] \;=\; \mathtt{jacrev}\!\bigl(u_\theta(\cdot,z)\bigr)\!(x_i)\,,
\]
$\text{per-iter cost: } \mathrm{latent\_dim}\times r\text{-eval} \;\to\; 2\times r\text{-eval}$;
measured fine-stage Jacobian time $54$\,ms $\to$ $17$\,ms.

\paragraph{\texttt{ModulatedSIREN} decoder (reference).}
Inherited from Session 1 and unchanged today, but worth a paragraph
because the linear-affine decoder below is best understood as its
contrast. The decoder is a SIREN multi-layer perceptron with sine
activations, $L{=}5$ layers and hidden width $h{=}384$, taking a 2-D
coordinate $x$ and returning a scalar $u(x)$. A small modulator MLP
(\texttt{mod\_in} $\to$ ReLU $\to$ \texttt{mod\_out}) maps the latent
$z$ to a stack of per-layer bias offsets $\beta_l \in \mathbb{R}^h$, so
each sine layer is
\[
h_{l+1} \;=\; \sin\!\bigl(\omega \cdot (W_l \, h_l + b_l + \beta_l(z))\bigr).
\]
The dependence on $z$ enters \emph{inside} the sines, so the map is
nonlinear in $z$ and $\partial u / \partial z$ varies with $z$ ---
which is why every GN step needs a fresh Jacobian.

\paragraph{\texttt{AffineZDecoder} (the failed first attempt).}
First try at making the Jacobian cheap: write
$u(x;z) = \Phi(x) \cdot A(z) + b(x)$ where $\Phi$ is a fixed SIREN
feature trunk (same depth/width as ModulatedSIREN), $A(z)$ is a small
ReLU MLP \texttt{latent} $\to$ \texttt{hidden\_dim}, and $b(x)$ is a
learned scalar mean field. $\Phi(x_{\text{eq}})$ can be precomputed
once per solve, so each GN iteration only needs $A(z)$ and
$\mathtt{jacrev}(A)(z)$ --- both tiny MLP calls, total wall $\sim$1\,ms
per iter. Per-sample reconstruction hit $\sim$1e-2 in training, but
ROM rel-L$^2$ stayed at $0.2$--$0.8$: $A$ is nonlinear, so the
Jacobian $\partial u / \partial z = \Phi(x) \cdot \mathrm{d}A/\mathrm{d}z$
is valid only in a small basin around $z{=}0$, and cold-start GN
walked out of that basin before converging. Widening hidden\_dim,
bumping latent to 32, and dropping the anchor loss none of them
escaped the basin --- the problem is structural, not a tuning issue.

\paragraph{\texttt{LinearAffineZDecoder} (the fix).}
Replace $A(z)$ with a single linear map $A(z) = W_A \, z + b_A$. Then
$u$ becomes linear in $z$ with $z$-independent coefficients:
\[
u(x; z) \;=\; \underbrace{\bigl(\Phi(x)\,W_A\bigr)}_{V(x)} \, z \;+\; \Phi(x)\,b_A \;+\; b(x),
\qquad \frac{\partial u}{\partial z} \;=\; V(x) \;\;\text{(constant in } z\text{)}.
\]
Structurally this is a Compact-POD decoder: $V(x)$ is the reduced
basis, $z$ is the coefficient vector, $b(x)$ is the mean field. The
difference vs classical CP is that $V$ is learned end-to-end with the
ViT encoder (rather than computed by SVD on snapshots) and $b$ is
learned (rather than the training-set mean). Because the Jacobian is
constant, the NM-ROM residual (nonlinear in $u$, but linear in $z$
through this decoder) reduces to one linear least-squares step per GN
iter; in practice GN converges in 2 iters. Paired with
\texttt{AffineNMROMSolver}, which precomputes $V(x_{\text{eq}})$ and
$b(x_{\text{eq}})$ once per solve. Trained with a Laplacian residual
loss $\lambda_{\text{lap}} \cdot \|K \cdot \mathrm{decode}(z) - F\|^2$
evaluated at random interior 5-point stencils per batch; without it
the encoder finds a good $z$ but the ROM-residual minimum is in a
different direction (rel-L$^2$ 0.25); with it the encoder-vs-ROM $\|z\|$
agreement is within 1\% and ROM rel-L$^2$ drops to $\sim$1e-2.

\paragraph{Jacobian behavior at a glance.}

\begin{center}
\begin{tabular}{l l l r}
  \textbf{Decoder} & \textbf{$\partial u / \partial z$} & \textbf{Per-iter Jacobian cost} & \textbf{Cold-start ROM} \\
  \hline
  ModulatedSIREN (jacfwd)       & nonlinear in $z$ & $\mathrm{latent\_dim} \times r$-eval & 2.88e-3 (with composite recipe) \\
  ModulatedSIREN (vmap jacrev)  & nonlinear in $z$ & $\sim 2 \times r$-eval              & 2.88e-3 (3.2$\times$ faster Jac)    \\
  AffineZDecoder                & nonlinear in $z$ & tiny MLP $+$ jacrev                  & 0.2--0.8 (basin trap)               \\
  LinearAffineZDecoder          & \emph{constant} $V(x)$ & one precompute, then free       & 9.81e-3 (with $\lambda_{\text{lap}}$) \\
\end{tabular}
\end{center}

